% ЗГЕНЕРОВАНО scripts/build_format_v2_chapter.py — не редагувати вручну.
% Джерело: editorial/cards/13.*.json, format-v2/rozdil-13.json, paraleli.json.

\naukchapter{13}{7}{Глава 13. Я сиджу, як мені завгодно}
\naukintro{Тринадцята глава — сім віршів і один рефрен, який стоїть у кінці кожного: я сиджу, як мені завгодно. Слово сиджу тут не про позу: глава саме й показує, що поза нічого не додає (13.5). Рефрен тримає іншу думку — що стан того, хто пізнав, не залежить ні від дії, ні від недіяння, і саме тому його не можна ані здобути, ані втратити.}
\vsec{13.1}
\vdeva{अकिंचनभवं स्वास्थ्यं कौपीनत्वे ऽपि दुर्लभम् । त्यागादाने विहायास्माद् अहम् आसे यथासुखम् ॥}
\viast{akiṃcanabhavaṃ svāsthyaṃ kaupīnatve 'pi durlabham | tyāgādāne vihāyāsmād aham āse yathāsukham ||}
\vtranslit{акімчанабгавам свастхям каупінатве ’пі дурлабгам | тяґадане віхаясмад ахам асе ятхасукхам ||}
\vperek{Стан неволодіння — це бути в собі; навіть у жебрацькому поясі він рідкісний. Відкинувши віддавання й приймання, я сиджу, як мені завгодно.}
\vcom{Построкове зіставлення: редакторська чернетка, не рецензовано.}
\vcom{Рядок а}
\vdeva{अकिंचनभवं स्वास्थ्यं कौपीनत्वे ऽपि दुर्लभम्}
\viast{akiṃcanabhavaṃ svāsthyaṃ kaupīnatve 'pi durlabham}
\vtranslit{акімчанабгавам свастхям каупінатве ’пі дурлабгам}
\vperek{Стан неволодіння — це бути в собі; навіть у жебрацькому поясі він рідкісний.}
\vcom{akiṃcana-bhavaṃ називає стан неволодіння, а svāsthyaṃ — перебування у власному стані; це не медичне твердження про здоров’я. kaupīnatve 'pi durlabham підкреслює, що навіть зовнішня ознака аскези (стегнова пов’язка) не гарантує внутрішньої свободи.}
\vcom{Рядок б}
\vdeva{त्यागादाने विहायास्माद् अहम् आसे यथासुखम्}
\viast{tyāgādāne vihāyāsmād aham āse yathāsukham}
\vtranslit{тяґадане віхаясмад ахам асе ятхасукхам}
\vperek{Відкинувши віддавання й приймання, я сиджу, як мені завгодно.}
\vcom{tyāgādāne — пара «віддавання й приймання»; vihāya означає «відкинувши». Після неї aham āse yathāsukham — повторюваний рефрен глави («я сиджу, як мені завгодно»). Думка стосується ставлення до володіння, а не буквальної вимоги сидіти.}
\vcom{Найважливіше тут — svāsthya. Слово означає буквально «перебування в собі» і в сучасній мові значить «здоров'я»; у вірші воно вжите в першому сенсі, і це не метафора, а визначення: бути здоровим — значить бути в собі. Друга точність — kaupīna: жебрацький пояс названий не як приклад аскези, а як те, у чому svāsthya все одно рідкісне. Тобто зовнішня ознака відречення нічого не гарантує. Пара tyāga / ādāna («віддавання» / «приймання») — та сама, що heya / upādeya у 7.5, і тут її теж відкидають цілою, а не по частинах. РЕФРЕН ГЛАВИ. Шість із семи віршів (13.1, 13.3–13.7) закінчуються однаково: aham āse yathāsukham («я сиджу — як мені завгодно»). ās — не «бути», а саме СИДІТИ, перебувати на місці; yathāsukham — «як зручно, за бажанням». Разом це формула, яка скасовує будь-яку поставу як вимогу: не треба ні стояти, ні лежати правильно. 13.2 — єдиний вірш глави без цього рефрену, і його треба читати окремо.}
\vpar{Українська паралель (наша інтерпретація)}{Григорій Сковорода, «Асхань. Симфонія, наречена книга Асхань, про пізнання самого себе».}{Наша паралель: у Сковороди бути в собі теж означає не мати нічого свого — і саме тому цей стан не залежить від достатку. Наше порівняння — типологічне.}
\vfig{ilustratsii/13-1-neti.pdf}{Віддавання й приймання відкинуті: сиджу, як завгодно}{Схема наша, не історична ілюстрація.}
\vfig{ilustratsii-naratyvni/13-1-sydyt-iak-zavhodno.pdf}{Людина сидить під деревом, а жебрацький пояс і чаша лежать осторонь}{Наша ілюстрація до вірша 13.1 — не історичне зображення.}

\vsec{13.2}
\vdeva{कुत्रापि खेदः कायस्य जिह्वा कुत्रापि खिद्यते । मनः कुत्रापि तत् त्यक्त्वा पुरुषार्थे स्थितः सुखम् ॥}
\viast{kutrāpi khedaḥ kāyasya jihvā kutrāpi khidyate | manaḥ kutrāpi tat tyaktvā puruṣārthe sthitaḥ sukham ||}
\vtranslit{кутрапі кхедах каяся джіхва кутрапі кхідяте | манах кутрапі тат тяктва пурушартхе стхітах сукхам ||}
\vperek{Десь втомлюється тіло, десь втомлюється язик, десь розум — відкинувши це, я стою в меті людини, щасливо.}
\vcom{Построкове зіставлення: редакторська чернетка, не рецензовано.}
\vcom{Рядок а}
\vdeva{कुत्रापि खेदः कायस्य जिह्वा कुत्रापि खिद्यते}
\viast{kutrāpi khedaḥ kāyasya jihvā kutrāpi khidyate}
\vtranslit{кутрапі кхедах каяся джіхва кутрапі кхідяте}
\vperek{Десь втомлюється тіло, десь втомлюється язик,}
\vcom{kutrāpi повторено двічі («десь»), зіставляючи khedaḥ kāyasya («втому тіла») та jihvā ... khidyate («втомлюється язик», тобто орган мовлення). Перший піврядок не містить ще висновку щодо того, що слід відкинути.}
\vcom{Рядок б}
\vdeva{मनः कुत्रापि तत् त्यक्त्वा पुरुषार्थे स्थितः सुखम्}
\viast{manaḥ kutrāpi tat tyaktvā puruṣārthe sthitaḥ sukham}
\vtranslit{манах кутрапі тат тяктва пурушартхе стхітах сукхам}
\vperek{десь розум — відкинувши це, я стою в меті людини, щасливо.}
\vcom{manaḥ kutrāpi додає третю сферу — розум; tat tyaktvā («відкинувши це») відсилає до названих утом. puruṣārthe sthitaḥ sukham — перебування в людській меті або її здійсненні, щасливо; це окрема думка, і вірш 13.2 не закінчується звичайним рефреном aham āse yathāsukham.}
\vcom{ЄДИНИЙ вірш глави без рефрену aham āse yathāsukham, і це не випадково: тут ідеться не про те, як сидіти, а про те, ЩО втомлюється. Три рівні названі поспіль — тіло, язик, розум — і кожен втомлюється «десь», тобто не завжди й не однаково. Слово kheda («втома») м'якше за duḥkha («страждання») з глави 8: це не біль, а зношування. puruṣārtha («мета людини») вжите в однині, хоч традиційно їх чотири, — тобто йдеться не про вибір між ними, а про те, що всі вони вже не потрібні.}
\vpar{Українська паралель (наша інтерпретація)}{Олександр Довженко, «Щоденник» (1941–1956).}{Наша паралель: у Довженка втома тіла, мови й думки теж описується окремо, і саме розрізнення цих трьох втом показує, що людина не зводиться до жодної з них. Наше порівняння — за будовою думки.}
\vfig{ilustratsii/13-2-neti.pdf}{Три втоми відкинуті: лишається мета людини}{Схема наша, не історична ілюстрація.}

\vsec{13.3}
\vdeva{कृतं किम् अपि नैव स्याद् इति संचिन्त्य तत्त्वतः । यदा यत् कर्तुम् आयाति तत् कृत्वासे यथासुखम् ॥}
\viast{kṛtaṃ kim api naiva syād iti saṃcintya tattvataḥ | yadā yat kartum āyāti tat kṛtvāse yathāsukham ||}
\vtranslit{крітам кім апі найва сяд іті самчінтя таттватах | яда ят картум аяті тат крітвасе ятхасукхам ||}
\vperek{Подумавши по-справжньому: «зробленого не буде нічого» — що б не випало зробити, зробивши те, я сиджу, як мені завгодно.}
\vcom{Построкове зіставлення: редакторська чернетка, не рецензовано.}
\vcom{Рядок а}
\vdeva{कृतं किम् अपि नैव स्याद् इति संचिन्त्य तत्त्वतः}
\viast{kṛtaṃ kim api naiva syād iti saṃcintya tattvataḥ}
\vtranslit{крітам кім апі найва сяд іті самчінтя таттватах}
\vperek{Подумавши по-справжньому: «зробленого не буде нічого» —}
\vcom{kṛtaṃ kim api naiva syād задає заперечення результату дії; iti saṃcintya tattvataḥ означає «так розваживши по суті». Оптатив syāt не варто без пояснення підмінювати твердженням, що жодна дія фізично не відбувається.}
\vcom{Рядок б}
\vdeva{यदा यत् कर्तुम् आयाति तत् कृत्वासे यथासुखम्}
\viast{yadā yat kartum āyāti tat kṛtvāse yathāsukham}
\vtranslit{яда ят картум аяті тат крітвасе ятхасукхам}
\vperek{що б не випало зробити, зробивши те, я сиджу, як мені завгодно.}
\vcom{yadā yat kartum āyāti описує те, що випадає зробити; tat kṛtvāse поєднує «зробивши те» з «сиджу». Рефрен yathāsukham збережено у чинному перекладі: відсутність привласнення результату не заперечує самого виконання дії.}
\vcom{Парадокс у тому, що дія не скасовується: yadā yat kartum āyāti («коли що випадає зробити») — дієслово āyāti означає саме те, що приходить само, без запрошення. Отже, діють, але без того, хто міг би зарахувати її собі. kṛtvāse — злитна форма kṛtvā + āse, і в ній дві дії стоять впритул: зробивши — сиджу. Умова naiva syāt («нехай не буде») стоїть у optativus, тобто це не констатація, а спосіб бачення, прийнятий заздалегідь. РЕФРЕН ГЛАВИ. Шість із семи віршів (13.1, 13.3–13.7) закінчуються однаково: aham āse yathāsukham («я сиджу — як мені завгодно»). ās — не «бути», а саме СИДІТИ, перебувати на місці; yathāsukham — «як зручно, за бажанням». Разом це формула, яка скасовує будь-яку поставу як вимогу: не треба ні стояти, ні лежати правильно. 13.2 — єдиний вірш глави без цього рефрену, і його треба читати окремо.}
\vpar{Українська паралель (ПОТРЕБУЄ ПЕРЕВІРКИ)}{українська філософська лірика, не визначено.}{ПОТРЕБУЄ ПЕРЕВІРКИ: думка, що зробленого не буде нічого, близька до української філософської лірики, але конкретного твору ми не звірили.}
\vfig{ilustratsii/13-3-space.pdf}{Небо без хмар: зробленого не буде нічого}{Схема наша, не історична ілюстрація.}

\vsec{13.4}
\vdeva{कर्मनैष्कर्म्यनिर्बन्धभावा देहस्थयोगिनः । संयोगायोगविरहाद् अहम् आसे यथासुखम् ॥}
\viast{karmanaiṣkarmyanirbandhabhāvā dehasthayoginaḥ | saṃyogāyogavirahād aham āse yathāsukham ||}
\vtranslit{карманайшкарм'янірбандгабгава дехастхайоґінах | самйоґайоґавірахад ахам асе ятхасукхам ||}
\vperek{У йога, що перебуває в тілі, є стани, до яких його змушують дія й недіяння; та через відсутність з'єднання й роз'єднання я сиджу, як мені завгодно.}
\vcom{Построкове зіставлення: редакторська чернетка, не рецензовано.}
\vcom{Рядок а}
\vdeva{कर्मनैष्कर्म्यनिर्बन्धभावा देहस्थयोगिनः}
\viast{karmanaiṣkarmyanirbandhabhāvā dehasthayoginaḥ}
\vtranslit{карманайшкарм'янірбандгабгава дехастхайоґінах}
\vperek{У йога, що перебуває в тілі, є стани, до яких його змушують дія й недіяння;}
\vcom{Складне слово karmanaiṣkarmyanirbandhabhāvā зіставляє дію (karma), недіяння (naiṣkarmya) й наполягання або примус (nirbandha); dehasthayoginaḥ вказує на йога, який перебуває в тілі. Це твердження про напругу між двома способами поведінки, а не доказ певної історичної школи.}
\vcom{Рядок б}
\vdeva{संयोगायोगविरहाद् अहम् आसे यथासुखम्}
\viast{saṃyogāyogavirahād aham āse yathāsukham}
\vtranslit{самйоґайоґавірахад ахам асе ятхасукхам}
\vperek{та через відсутність з'єднання й роз'єднання я сиджу, як мені завгодно.}
\vcom{saṃyogāyogavirahād — аблативна конструкція з ідеєю відсутності з'єднання й нез'єднання. Далі aham āse yathāsukham повертає рефрен особистого перебування; не слід перетворювати це на припис відмовитися від щоденних обов’язків.}
\vcom{Слово nirbandha («змушування») — ключове: доки є тіло, є й примус то до дії, то до недіяння, і жоден вибір не вільний. saṃyoga / ayoga («з'єднання» / «роз'єднання») — пара з мови йоги, де з'єднання зазвичай означає зв'язок з об'єктом; тут відкидається сама можливість такого зв'язку. Зверни увагу на порядок: спершу визнано примус (він є), потім названо те, чого немає (з'єднання). Вірш не обіцяє свободи від тіла — лише від рахунку. РЕФРЕН ГЛАВИ. Шість із семи віршів (13.1, 13.3–13.7) закінчуються однаково: aham āse yathāsukham («я сиджу — як мені завгодно»). ās — не «бути», а саме СИДІТИ, перебувати на місці; yathāsukham — «як зручно, за бажанням». Разом це формула, яка скасовує будь-яку поставу як вимогу: не треба ні стояти, ні лежати правильно. 13.2 — єдиний вірш глави без цього рефрену, і його треба читати окремо.}
\vpar{Українська паралель (ПОТРЕБУЄ ПЕРЕВІРКИ)}{українська релігійна проза, не визначено.}{ПОТРЕБУЄ ПЕРЕВІРКИ: розрізнення з'єднання й роз'єднання як того, що змушує людину до станів, найближче до української релігійної прози, але назвати твір без звірки ми не можемо.}
\vfig{ilustratsii/13-4-knot.pdf}{Вузол дії й недіяння розвʼязано}{Схема наша, не історична ілюстрація.}

\vsec{13.5}
\vdeva{अर्थानर्थौ न मे स्थित्या गत्या न शयनेन वा । तिष्ठन् गच्छन् स्वपन् तस्माद् अहम् आसे यथासुखम् ॥}
\viast{arthānarthau na me sthityā gatyā na śayanena vā | tiṣṭhan gacchan svapan tasmād aham āse yathāsukham ||}
\vtranslit{артханартхау на ме стхітя ґатя на шаянена ва | тіштхан ґаччан свапан тасмад ахам асе ятхасукхам ||}
\vperek{Ні користі, ні шкоди немає мені від стояння, ходіння чи лежання — тому, стоячи, ходячи чи сплячи, я сиджу, як мені завгодно.}
\vcom{Построкове зіставлення: редакторська чернетка, не рецензовано.}
\vcom{Рядок а}
\vdeva{अर्थानर्थौ न मे स्थित्या गत्या न शयनेन वा}
\viast{arthānarthau na me sthityā gatyā na śayanena vā}
\vtranslit{артханартхау на ме стхітя ґатя на шаянена ва}
\vperek{Ні користі, ні шкоди немає мені від стояння, ходіння чи лежання —}
\vcom{arthānarthau — двоїна «користь і шкода»; sthityā, gatyā, śayanena — орудні форми «стоянням, рухом, лежанням». Це оцінка способів тілесного перебування, а не заборона ходити чи лежати.}
\vcom{Рядок б}
\vdeva{तिष्ठन् गच्छन् स्वपन् तस्माद् अहम् आसे यथासुखम्}
\viast{tiṣṭhan gacchan svapan tasmād aham āse yathāsukham}
\vtranslit{тіштхан ґаччан свапан тасмад ахам асе ятхасукхам}
\vperek{тому, стоячи, ходячи чи сплячи, я сиджу, як мені завгодно.}
\vcom{tiṣṭhan, gacchan, svapan — три дієприкметники «стоячи, йдучи, сплячи»; tasmād — «тому». Слідом aham āse yathāsukham зберігає словесний парадокс «сиджу» попри зміну постави; рефрен уособлює сталість внутрішнього ставлення, а не конкретну фізичну позу.}
\vcom{Найгостріший парадокс глави, і він у самій граматиці: три дієприкметники (tiṣṭhan, gacchan, svapan — «стоячи, ходячи, сплячи») стоять поруч із дієсловом āse («сиджу»). Формально це суперечність, і вона навмисна: ās тут уже не про поставу тіла, а про перебування взагалі. Пара artha / anartha («користь» / «шкода») — та сама, що в 7.2 vṛddhi / kṣati, і знову відкидається цілою: не «немає шкоди», а немає ні того, ні того. РЕФРЕН ГЛАВИ. Шість із семи віршів (13.1, 13.3–13.7) закінчуються однаково: aham āse yathāsukham («я сиджу — як мені завгодно»). ās — не «бути», а саме СИДІТИ, перебувати на місці; yathāsukham — «як зручно, за бажанням». Разом це формула, яка скасовує будь-яку поставу як вимогу: не треба ні стояти, ні лежати правильно. 13.2 — єдиний вірш глави без цього рефрену, і його треба читати окремо.}
\vpar{Українська паралель (наша інтерпретація)}{Василь Стефаник, «Камінний хрест» (1900).}{Наша паралель: у Стефаника тілесне положення людини теж не міняє того, що з нею сталося всередині. Наше порівняння — типологічне.}
\vfig{ilustratsii/13-5-rest.pdf}{Людина під зайшлим сонцем: поза не додає ні користі, ні шкоди}{Схема наша, не історична ілюстрація.}

\vsec{13.6}
\vdeva{स्वपतो नास्ति मे हानिः सिद्धिर् यत्नवतो न वा । नाशोल्लासौ विहायास्माद् अहम् आसे यथासुखम् ॥}
\viast{svapato nāsti me hāniḥ siddhir yatnavato na vā | nāśollāsau vihāyāsmād aham āse yathāsukham ||}
\vtranslit{свапато насті ме ханіх сіддгір ятнавато на ва | нашолласау віхаясмад ахам асе ятхасукхам ||}
\vperek{У сплячого немає мені втрати, як і немає досягнення в того, хто докладає зусиль; відкинувши загибель і піднесення, я сиджу, як мені завгодно.}
\vcom{Построкове зіставлення: редакторська чернетка, не рецензовано.}
\vcom{Рядок а}
\vdeva{स्वपतो नास्ति मे हानिः सिद्धिर् यत्नवतो न वा}
\viast{svapato nāsti me hāniḥ siddhir yatnavato na vā}
\vtranslit{свапато насті ме ханіх сіддгір ятнавато на ва}
\vperek{У сплячого немає мені втрати, як і немає досягнення в того, хто докладає зусиль;}
\vcom{svapataḥ (у сандгі svapato) — родовий однини дієприкметника «той, хто спить»; na asti — «немає»; me — «мені / в мене»; hāniḥ — «втрата». siddhiḥ (у сандгі siddhir) — «успіх, досягнення»; yatnavataḥ (у сандгі yatnavato) — «того, хто докладає зусиль»; na vā завершує заперечення щодо досягнення. Це не наказ спати або відмовитися від праці: рядок зіставляє дві оцінки — втрату й здобуток.}
\vcom{Рядок б}
\vdeva{नाशोल्लासौ विहायास्माद् अहम् आसे यथासुखम्}
\viast{nāśollāsau vihāyāsmād aham āse yathāsukham}
\vtranslit{нашолласау віхаясмад ахам асе ятхасукхам}
\vperek{відкинувши загибель і піднесення, я сиджу, як мені завгодно.}
\vcom{nāśa-ullāsau — двоїна, «загибель і піднесення»; vihāya — дієприслівникова форма «відкинувши». Форма vihāyāsmād передає злиття vihāya + asmāt («від цього / тому»); точний синтаксичний зв'язок варто окремо перевірити за друкованими джерелами. aham āse yathāsukham — «я сиджу, як мені завгодно». Рефрен глави не скасовує потреби уважно розглядати межу слів і сандгі.}
\vcom{Вірш будується на запереченні двох протилежних чеснот: сон не шкодить, зусилля не дає. Тим самим знімається сама вісь, на якій тримається будь-яка практика, — «старанність». ullāsa («піднесення») — слово з відтінком радості, не просто зростання; у парі з nāśa («загибель») воно описує не події, а те, як людина їх переживає. Порівняй 3.9, де похвала й мука були описані однією граматичною формою: тут той самий хід, але на рівні оцінки. РЕФРЕН ГЛАВИ. Шість із семи віршів (13.1, 13.3–13.7) закінчуються однаково: aham āse yathāsukham («я сиджу — як мені завгодно»). ās — не «бути», а саме СИДІТИ, перебувати на місці; yathāsukham — «як зручно, за бажанням». Разом це формула, яка скасовує будь-яку поставу як вимогу: не треба ні стояти, ні лежати правильно. 13.2 — єдиний вірш глави без цього рефрену, і його треба читати окремо.}
\vpar{Українська паралель (наша інтерпретація)}{Ліна Костенко, «Берестечко» (1999).}{Наша паралель: у Костенко теж є стан, у якому ні сон, ні зусилля не дають ні втрати, ні досягнення. Наше порівняння — за будовою думки.}
\vfig{ilustratsii/13-6-balance.pdf}{Дві шальки: втрата й досягнення — обидві відкинуті}{Схема наша, не історична ілюстрація.}

\vsec{13.7}
\vdeva{सुखादिरूपा नियमं भावेष्व् आलोक्य भूरिशः । शुभाशुभे विहायास्माद् अहम् आसे यथासुखम् ॥}
\viast{sukhādirūpā niyamaṃ bhāveṣv ālokya bhūriśaḥ | śubhāśubhe vihāyāsmād aham āse yathāsukham ||}
\vtranslit{сукхадірупа ніямам бгавешв алокя бгурішах | шубгашубге віхаясмад ахам асе ятхасукхам ||}
\vperek{Побачивши не раз у всіх станах те правило, що має вигляд задоволення, — відкинувши добре й погане, я сиджу, як мені завгодно.}
\vcom{Построкове зіставлення: редакторська чернетка, не рецензовано.}
\vcom{Рядок а}
\vdeva{सुखादिरूपा नियमं भावेष्व् आलोक्य भूरिशः}
\viast{sukhādirūpā niyamaṃ bhāveṣv ālokya bhūriśaḥ}
\vtranslit{сукхадірупа ніямам бгавешв алокя бгурішах}
\vperek{Побачивши не раз у всіх станах те правило, що має вигляд задоволення, —}
\vcom{bhāveṣv — місцевий множини «у станах»; ālokya — дієприслівникова форма «побачивши»; bhūriśaḥ — «багаторазово». Чинна IAST-картка має niyamaṃ («правило»), що відображено в її українському перекладі; SARIT має aniyamaṃ («відсутність правила»). Це змістове різночитання не розв'язано, а sukhādirūpā також має проблемне граматичне узгодження. Построковість фіксує читання картки, не видає його за доведений оригінал.}
\vcom{Рядок б}
\vdeva{शुभाशुभे विहायास्माद् अहम् आसे यथासुखम्}
\viast{śubhāśubhe vihāyāsmād aham āse yathāsukham}
\vtranslit{шубгашубге віхаясмад ахам асе ятхасукхам}
\vperek{відкинувши добре й погане, я сиджу, як мені завгодно.}
\vcom{śubhāśubhe — двоїна «добре й погане»; vihāya — «відкинувши»; у vihāyāsmād видно сандгі з asmāt («із цього / тому»). aham āse — «я сиджу», а yathāsukham — «як мені зручно». Це повторюваний фінал глави 13, але він не вирішує різночитання niyamaṃ / aniyamaṃ у першому піврядку.}
\vcom{Найважче місце глави, і не через сенс, а через текст. Свідки розходяться в одному слові: SARIT подає aniyamaṃ («безладдя, відсутність правила»), GRETIL — niyamaṃ («правило»). Сенс при цьому стає ПРОТИЛЕЖНИМ: побачити, що все, що має вигляд задоволення, насправді безладдя — або що воно підлягає правилу. Крім того, форма sukhādirūpā (жіночий рід, однина) не узгоджується ні з niyamaṃ/aniyamaṃ (середній рід), ні з bhāveṣu (множина, чоловічий/середній). Найімовірніше, обидва свідки успадкували зіпсоване читання, і правильним було sukhādirūpāṇi (середній рід, множина) — «речі, чия подоба — задоволення тощо». У перекладі взято читання SARIT, але це ВИБІР, а не розв'язання: потрібен третій свідок. РЕФРЕН ГЛАВИ. Шість із семи віршів (13.1, 13.3–13.7) закінчуються однаково: aham āse yathāsukham («я сиджу — як мені завгодно»). ās — не «бути», а саме СИДІТИ, перебувати на місці; yathāsukham — «як зручно, за бажанням». Разом це формула, яка скасовує будь-яку поставу як вимогу: не треба ні стояти, ні лежати правильно. 13.2 — єдиний вірш глави без цього рефрену, і його треба читати окремо.}
\vpar{Українська паралель (ПОТРЕБУЄ ПЕРЕВІРКИ)}{українська барокова поезія, не визначено.}{ПОТРЕБУЄ ПЕРЕВІРКИ: правило, що має вигляд задоволення, близьке до барокового мотиву оманливої втіхи, але конкретного тексту ми не перевірили.}
\vfig{ilustratsii/13-7-illusion.pdf}{Серпанок задоволення: під ним — добре й погане}{Схема наша, не історична ілюстрація.}

